\chapter{Carpal Tunnel Release}

\section*{What Is This Surgery?}
Carpal tunnel release is a frequently performed surgical procedure aimed at alleviating the debilitating symptoms of carpal tunnel syndrome (CTS). This condition arises from the compression of the median nerve as it traverses the carpal tunnel, a confined osteofibrous passageway located at the wrist. The surgery involves the precise division of the transverse carpal ligament, which forms the roof of this tunnel, thereby effectively increasing the available space for the median nerve and the nine flexor tendons. By decompressing the nerve, the procedure seeks to restore normal nerve function, mitigate pain, numbness, and tingling sensations in the hand and fingers, and prevent the progression of nerve damage or muscle atrophy. It represents a highly effective and definitive intervention for patients whose symptoms are refractory to conservative management or who present with objective signs of severe median nerve compromise.

\section*{Alternative Names}
\begin{itemize}
  \item Carpal Tunnel Decompression
  \item Median Nerve Decompression
  \item Transverse Carpal Ligament Release
  \item Open Carpal Tunnel Release (OCTR)
  \item Endoscopic Carpal Tunnel Release (ECTR)
  \item Wrist Decompression Surgery for CTS
\end{itemize}

\section*{Relevant Surgical Anatomy}
The carpal tunnel is a critical anatomical structure at the wrist, forming a narrow passageway for vital neurovascular elements. Its boundaries are defined by the carpal bones forming a concave arch posteriorly and laterally, and the strong, unyielding transverse carpal ligament (also known as the flexor retinaculum) anteriorly, which spans from the scaphoid and trapezium radially to the pisiform and hook of the hamate ulnarly. Within this confined space, the median nerve is situated most superficially, just beneath the transverse carpal ligament, making it vulnerable to compression. Deep to the median nerve lie the nine flexor tendons: four tendons of the flexor digitorum superficialis, four tendons of the flexor digitorum profundus, and the single tendon of the flexor pollicis longus, all enveloped by synovial sheaths.

The median nerve provides motor innervation to the thenar muscles (abductor pollicis brevis, opponens pollicis, superficial head of flexor pollicis brevis) via its recurrent motor branch, which typically arises from the radial side of the nerve distal to the transverse carpal ligament, but can have variable anatomy, sometimes piercing or running through the ligament. Sensory innervation is provided to the palmar aspect of the thumb, index finger, middle finger, and the radial half of the ring finger, as well as the corresponding palm. The palmar cutaneous branch of the median nerve, which provides sensation to the proximal palm, typically arises proximal to the carpal tunnel and travels superficial to the transverse carpal ligament, thus usually spared during surgery. Important surrounding structures to consider during surgery include the ulnar nerve and artery (located ulnar to the carpal tunnel, within Guyon's canal), and the superficial palmar arterial arch, which lies distal to the transverse carpal ligament. Surgical landmarks include the thenar crease, the distal wrist crease, and the palpable hook of the hamate, which guides the ulnar extent of the transverse carpal ligament.

\section{Surgical Principle \& Rationale}
The fundamental surgical principle of carpal tunnel release is to irreversibly enlarge the volume of the carpal tunnel, thereby decompressing the median nerve. This is achieved by surgically transecting the transverse carpal ligament, which acts as the unyielding roof of the tunnel. In carpal tunnel syndrome, any increase in pressure within this fixed-volume space, whether due to tenosynovitis, edema, trauma, or systemic conditions, directly compromises the median nerve's microcirculation and axonal transport, leading to ischemia, demyelination, and ultimately axonal degeneration.

The rationale for surgical intervention is to promptly and effectively relieve this extrinsic compression. By dividing the transverse carpal ligament, the pressure on the median nerve is immediately reduced, allowing for improved blood flow, resolution of nerve edema, and restoration of normal nerve conduction. The ligament typically does not re-form in its original constricted state; instead, the gap created by its division fills with less restrictive scar tissue, maintaining the enlarged tunnel volume over the long term. The primary technical goals are the complete and safe division of the entire transverse carpal ligament from its proximal to distal attachments, while meticulously protecting the median nerve, its crucial recurrent motor branch, and other adjacent neurovascular structures from iatrogenic injury.

\section{History \& Surgical Evolution}
The clinical entity of median nerve compression at the wrist was first recognized by James Paget in 1854, who described median nerve symptoms following a distal radius fracture. However, it was Pierre Marie and Charles Foix in 1913 who more precisely characterized carpal tunnel syndrome and its association with median nerve compression within the carpal tunnel. Early surgical attempts in the first half of the 20th century were often limited to neurolysis or excision of specific space-occupying lesions.

A pivotal moment in the history of carpal tunnel surgery occurred in the 1950s when George S. Phalen, an American orthopedic surgeon, popularized the open surgical division of the transverse carpal ligament. Phalen's extensive work and publications established the efficacy of this procedure and standardized the open approach, making it the definitive treatment for CTS. The technique continued to evolve with the introduction of endoscopic carpal tunnel release (ECTR) in the 1980s by Okutsu and others. ECTR offered a less invasive alternative, utilizing small incisions and an endoscope to visualize and divide the ligament, promising potentially faster recovery, reduced postoperative pain, and less scar tenderness compared to the traditional open method. Both open and endoscopic techniques have undergone continuous refinement in instrumentation and surgical approaches, with ongoing research focusing on optimizing patient selection, minimizing complications, and enhancing long-term functional outcomes.

\section{Clinical Indications}
Carpal tunnel release is indicated for patients presenting with symptoms of median nerve compression that significantly impact their quality of life or show signs of progressive nerve damage.

\begin{itemize}
  \item Persistent or worsening symptoms of carpal tunnel syndrome, including numbness, tingling, pain, or burning sensations in the median nerve distribution (thumb, index, middle, and radial half of the ring finger), especially if nocturnal.
  \item Failure of conservative management, such as wrist splinting, activity modification, non-steroidal anti-inflammatory drugs (NSAIDs), and corticosteroid injections, after an adequate trial period, typically 3-6 months.
  \item Objective evidence of median nerve compression on electrodiagnostic studies (nerve conduction studies and electromyography), confirming the diagnosis and assessing the severity of nerve involvement.
  \item Presence of thenar muscle weakness or atrophy, indicating severe or long-standing median nerve compression and impending irreversible nerve damage.
  \item Constant or nocturnal pain that significantly disrupts sleep, work, or daily activities, despite non-surgical interventions.
  \item Progressive sensory loss or motor deficits, such as dropping objects or difficulty with fine motor tasks, despite non-surgical management.
  \item Acute carpal tunnel syndrome resulting from trauma (e.g., distal radius fracture with significant swelling), infection, or hemorrhage, requiring urgent decompression to prevent permanent nerve injury.
  \item Recurrence of symptoms after previous non-surgical treatment or incomplete initial surgical release.
\end{itemize}

\section{Clinical Positioning}
Carpal tunnel release is generally considered a primary definitive treatment for carpal tunnel syndrome when conservative measures have proven ineffective or when there is evidence of severe nerve compromise. It is typically not the initial intervention for newly diagnosed, mild cases, where non-surgical treatments such as nocturnal wrist splinting, activity modification, ergonomic adjustments, and corticosteroid injections are first-line therapies. However, in specific scenarios, such as patients presenting with significant motor weakness, thenar muscle atrophy, or severe, unremitting sensory deficits, surgical decompression may be considered a primary intervention to prevent irreversible nerve damage and optimize functional recovery. In most cases, it serves as a secondary or salvage procedure, offered after a thorough trial of conservative therapies has been exhausted without adequate symptom relief, or when symptoms recur after initial non-surgical success. The decision to proceed with surgery is always individualized, balancing the severity and duration of symptoms, objective electrodiagnostic findings, and patient preferences and expectations.

\section{Surgical Technique \& Procedure}
Carpal tunnel release is typically performed on an outpatient basis under local anesthesia, often supplemented with intravenous sedation for patient comfort. The patient's arm is positioned on a hand table, and a tourniquet is applied to the upper arm to ensure a bloodless surgical field, enhancing visualization and precision.

For the \textbf{Open Carpal Tunnel Release (OCTR)}:
\begin{itemize}
  \item \textbf{Anesthesia \& Positioning:} Local anesthetic (e.g., lidocaine with epinephrine) is infiltrated into the palm. The arm is supinated and secured on a hand table.
  \item \textbf{Incision:} A small longitudinal incision, typically 2-4 cm long, is made in the palm, usually along the ulnar border of the thenar crease or slightly ulnar to the midline of the ring finger.
  \item \textbf{Dissection:} The incision is carefully deepened through the skin and subcutaneous tissue, meticulously protecting superficial nerves (e.g., palmar cutaneous branch of the median nerve) and vessels.
  \item \textbf{Ligament Identification:} The palmar fascia is incised, exposing the underlying transverse carpal ligament, which appears as a thick, glistening white fibrous band.
  \item \textbf{Ligament Division:} Using a scalpel or fine scissors, the transverse carpal ligament is meticulously divided longitudinally from its distal attachment (often near the superficial palmar arch) to its proximal attachment (proximal to the distal wrist crease), ensuring complete release. Extreme care is taken to protect the median nerve and its recurrent motor branch, which is often visualized emerging from the radial side of the median nerve.
  \item \textbf{Confirmation:} The surgeon visually confirms adequate decompression of the median nerve, which often bulges slightly once the pressure is relieved.
  \item \textbf{Closure:} The skin incision is closed with fine non-absorbable sutures or surgical staples. A sterile dressing is applied, often incorporating a bulky soft dressing or a light splint. No drains are typically used.
\end{itemize}

For the \textbf{Endoscopic Carpal Tunnel Release (ECTR)}:
\begin{itemize}
  \item \textbf{Anesthesia \& Positioning:} Similar to OCTR, local anesthesia with sedation is common.
  \item \textbf{Incision(s):} This technique involves one or two small incisions (portals), typically less than 1 cm each. A common approach uses a single distal wrist crease incision (single-portal technique) or a proximal palmar incision (two-portal technique).
  \item \textbf{Endoscope Insertion:} An endoscope (a small camera) is inserted into the carpal tunnel through the portal, allowing the surgeon to visualize the transverse carpal ligament and the median nerve on a monitor.
  \item \textbf{Instrument Insertion:} A specialized cutting instrument (e.g., a hooked knife or blade) is then introduced through the same or a separate portal, guided by endoscopic visualization.
  \item \textbf{Ligament Division:} Under direct endoscopic visualization, the transverse carpal ligament is carefully divided from distal to proximal, ensuring complete release while protecting the median nerve and surrounding structures.
  \item \textbf{Confirmation:} The surgeon confirms complete division of the ligament and decompression of the median nerve by observing the nerve's appearance and the gap created.
  \item \textbf{Closure:} The small incisions are closed with one or two sutures or adhesive strips, and a sterile dressing is applied. No drains are typically used.
\end{itemize}

\section*{Preoperative Workup \& Preparation}
Prior to carpal tunnel release, a comprehensive preoperative workup is essential to confirm the diagnosis, assess the patient's overall health, and minimize surgical risks.
\begin{itemize}
  \item \textbf{Diagnosis Confirmation:} Electrodiagnostic studies, including nerve conduction studies (NCS) and electromyography (EMG), are almost always performed to objectively confirm the diagnosis of carpal tunnel syndrome, localize the site of median nerve compression, and assess the severity of nerve involvement. This helps differentiate CTS from other neuropathies (e.g., cervical radiculopathy).
  \item \textbf{Medical Clearance:} Patients with significant medical comorbidities (e.g., cardiovascular disease, diabetes, pulmonary issues) may require medical clearance from their primary care physician or a specialist (e.g., cardiologist). This may involve blood tests (e.g., complete blood count, basic metabolic panel, coagulation profile), electrocardiogram (ECG), and chest X-ray, as deemed necessary by the anesthesia team.
  \item \textbf{Medication Review:} A thorough review of all current medications is crucial. Patients are typically instructed to discontinue blood-thinning medications (e.g., aspirin, NSAIDs, clopidogrel, warfarin, novel oral anticoagulants) for a specified period (usually 5-10 days) before surgery to minimize bleeding risk. Diabetic patients receive specific instructions on managing their blood glucose levels around the time of surgery.
  \item \textbf{NPO Status:} If intravenous sedation or general anesthesia is planned, patients must strictly adhere to "nil per os" (NPO) guidelines, typically refraining from food and drink for 6-8 hours prior to surgery to prevent aspiration.
  \item \textbf{Prophylactic Antibiotics:} Intravenous prophylactic antibiotics (e.g., cefazolin) are commonly administered within 60 minutes prior to the surgical incision to reduce the risk of surgical site infection.
  \item \textbf{Informed Consent:} A detailed discussion with the surgeon covers the specifics of the procedure (open vs. endoscopic), potential benefits, anticipated risks and complications, alternative treatments, and the expected postoperative course and recovery timeline. The patient must fully understand and agree to the procedure by providing informed consent.
\end{itemize}

\section*{Contraindications \& Precautions}
\textbf{Absolute Contraindications:}
\begin{itemize}
  \item Active infection in the surgical field (e.g., cellulitis, abscess) or systemic sepsis, which must be treated prior to elective surgery.
  \item Uncontrolled bleeding disorder or an inability to safely discontinue anticoagulation medications for the required preoperative period.
  \item Patient refusal or inability to provide informed consent due to cognitive impairment without a legal guardian.
  \item Severe, uncorrectable medical comorbidities that render the risks of anesthesia or surgery excessively high, outweighing the potential benefits.
\end{itemize}

\textbf{Relative Contraindications and Precautions:}
\begin{itemize}
  \item Mild symptoms that respond well to conservative management, making surgery unnecessary at that time.
  \item Coexisting peripheral neuropathy from other causes (e.g., severe diabetes, alcoholism, vitamin deficiencies) that may complicate diagnosis, mask the true extent of CTS, or limit the potential for functional recovery.
  \item Prior surgery in the same carpal tunnel, which may result in significant scar tissue, altered anatomy, and increased risk of nerve injury, potentially favoring an open approach over endoscopic.
  \item Significant anatomical variations or the presence of space-occupying lesions (e.g., large ganglion cyst, lipoma, tumor) within the carpal tunnel, which may necessitate an open approach for better visualization and excision.
  \item Severe thenar muscle atrophy with complete denervation on EMG, where the potential for significant motor functional recovery may be limited despite successful decompression.
  \item Pregnancy: While not an absolute contraindication, surgery is often deferred until after delivery due as pregnancy-induced CTS is often transient and resolves postpartum, and to minimize potential risks of anesthesia to the fetus.
  \item Inexperience of the surgeon with endoscopic techniques, which may increase the risk of complications compared to an open approach.
\end{itemize}

\section*{Complications \& Risks}
While carpal tunnel release is generally considered a safe and effective procedure, as with any surgical intervention, there are potential risks and complications that patients should be aware of.

\begin{itemize}
  \item \textbf{General Surgical Risks:}
  \begin{itemize}
    \item Bleeding: Minor bruising and swelling are common; significant hemorrhage requiring intervention is rare.
    \item Infection: Superficial wound infection (cellulitis) or, rarely, deep space infection (e.g., tenosynovitis, osteomyelitis) can occur.
    \item Anesthesia complications: Reactions to local anesthetic, nausea, vomiting, or more serious cardiac or respiratory events, though rare with regional anesthesia and sedation.
    \item Scarring: Formation of a visible scar, which may be tender or hypertrophic.
  \end{itemize}
  \item \textbf{Procedure-Specific Risks:}
  \begin{itemize}
    \item Nerve injury: This is the most significant specific risk. It can involve direct laceration or contusion of the median nerve itself, its recurrent motor branch (leading to permanent thenar muscle weakness or atrophy), or other sensory nerves (e.g., palmar cutaneous branch, ulnar nerve).
    \item Incomplete release: Failure to fully divide the entire transverse carpal ligament, leading to persistent or recurrent symptoms. This may necessitate revision surgery.
    \item Pillar pain: Pain and tenderness in the thenar and hypothenar eminences (the "pillars" of the carpal tunnel), often exacerbated by gripping or weight-bearing, which can persist for several months postoperatively.
    \item Reflex Sympathetic Dystrophy (RSD) / Complex Regional Pain Syndrome (CRPS): A rare but severe chronic pain condition affecting the hand and wrist, characterized by pain, swelling, stiffness, and skin changes.
    \item Tendon injury: Accidental laceration or damage to one of the flexor tendons, though rare with careful technique.
    \item Vascular injury: Damage to the superficial palmar arch or other small vessels, leading to hematoma formation.
    \item Recurrence of symptoms: Although uncommon (1-5\%), symptoms can return years after successful surgery, often due to new causes of compression, re-scarring, or incomplete initial release.
    \item Stiffness: Temporary stiffness of the fingers or wrist, usually managed with early mobilization and hand therapy.
    \item Bowstringing of flexor tendons: A rare complication where the flexor tendons may bowstring slightly after ligament release, potentially causing mild weakness or discomfort.
  \end{itemize}
\end{itemize}

\section*{Postoperative Recovery Timeline}
Immediately following carpal tunnel release, the patient's hand will be bandaged, often with a bulky soft dressing, and typically elevated to minimize swelling. Pain medication will be prescribed to manage discomfort. Most patients are discharged home within a few hours, with detailed instructions for wound care, activity restrictions, and warning signs. During the first 24-48 hours, it is crucial to keep the hand elevated above heart level as much as possible and to avoid strenuous activities, heavy gripping, or repetitive wrist movements. Patients are generally encouraged to gently move their fingers to prevent stiffness and promote circulation. The initial dressing is typically kept clean and dry.

Within the first 1-2 weeks, the sutures are usually removed during a follow-up visit with the surgeon or a hand therapist. At this stage, patients are often advised to begin light hand exercises to restore range of motion and gradually increase activity levels, while still avoiding heavy lifting, forceful gripping, or repetitive tasks. Many patients experience prompt relief of nocturnal pain and paresthesias within days to weeks. Long-term recovery involves a gradual return to full activity over several weeks to months, with continued improvement in strength and sensation. Scar remodeling and resolution of pillar pain can take several months, and some mild discomfort may persist for up to 6 months. Physical or occupational therapy, specifically hand therapy, may be recommended to optimize recovery, especially for patients with significant preoperative weakness, stiffness, or those requiring an accelerated return to work or specific activities. Full recovery of strength and sensation, particularly in cases with severe preoperative nerve damage, can take up to a year.

\section*{Surgical Findings \& Pathology}
During carpal tunnel release, the surgeon's primary finding is the transverse carpal ligament, which is typically thickened, taut, and compressing the underlying median nerve. The median nerve itself may appear flattened, discolored (pale or hyperemic), or edematous, indicating chronic compression. Upon division of the ligament, the nerve often visibly expands or "bows out," confirming successful decompression. The flexor tendons within the tunnel may show signs of tenosynovitis, appearing inflamed or thickened. The surgeon will document the appearance of the nerve, the completeness of the ligament release, and any other incidental findings (e.g., ganglion cysts, lipomas, or anatomical variations).

In the vast majority of carpal tunnel release procedures, no tissue is sent for formal pathological examination, as the transverse carpal ligament is simply divided in situ. However, if an unusual mass or lesion is encountered within the carpal tunnel (e.g., a ganglion cyst, lipoma, or suspected tumor), a biopsy or excision of that specific tissue would be performed. In such cases, the pathology report would detail the histological findings of the resected mass, confirming its benign or malignant nature and providing a definitive diagnosis. For routine carpal tunnel release, the "pathology" is the clinical and intraoperative observation of median nerve compression by a thickened transverse carpal ligament.

\section{Expected Postoperative Course}
A normal and successful postoperative course following carpal tunnel release is characterized by a progressive and sustained improvement in symptoms and hand function. Patients typically experience prompt relief of nocturnal pain and paresthesias (numbness and tingling) within the first few days to weeks after surgery. Daytime symptoms, such as numbness during activities, also gradually diminish. While some incisional pain and pillar pain are common in the initial weeks, these should steadily improve. The hand and wrist will gradually regain full range of motion, and grip and pinch strength will progressively increase over several months as the median nerve recovers and muscles reinnervate.

The surgical wound should heal cleanly without signs of infection, and scar tenderness will typically subside over time. Patients should be able to resume light daily activities within 1-2 weeks and gradually return to more strenuous activities, including work and sports, over 6-12 weeks, depending on the demands of their occupation and the severity of their preoperative condition. The ultimate goal is the restoration of functional use of the hand for daily activities, work, and hobbies, with minimal to no residual symptoms of median nerve compression. While complete resolution of all symptoms may take several months, the overall trajectory should be one of continuous improvement.

\section{Postoperative Care \& Follow-up}
Postoperative care for carpal tunnel release is crucial for optimal healing and functional recovery.
\begin{itemize}
  \item \textbf{Wound Care:} The initial bulky dressing is typically removed within 24-48 hours, and a lighter dressing is applied. Patients are instructed to keep the wound clean and dry, avoiding immersion in water until sutures are removed and the wound is fully closed.
  \item \textbf{Suture Removal:} The first follow-up visit usually occurs 10-14 days post-surgery for wound inspection and suture removal.
  \item \textbf{Clinical Reassessment:} Subsequent visits, typically at 4-6 weeks and 3-6 months, assess symptom resolution, hand function, grip strength, and overall patient satisfaction.
  \item \textbf{Physical/Occupational Therapy:} For patients with significant preoperative deficits, persistent stiffness, or those requiring accelerated recovery for work or specific activities, a referral to a certified hand therapist may be made. Therapists guide patients through a structured exercise program to improve range of motion, strength, dexterity, and provide scar management techniques.
  \item \textbf{Activity Resumption:} Patients are gradually cleared to resume normal activities, with specific guidance on avoiding heavy lifting, forceful gripping, or repetitive tasks for several weeks to months, depending on individual recovery and the demands of their activities.
  \item \textbf{Warning Signs:} Patients are instructed to contact their surgeon immediately if they experience signs of infection (e.g., increasing redness, warmth, swelling, pus drainage, fever), severe or worsening pain unresponsive to medication, new or worsening numbness or weakness, or excessive bleeding from the wound.
\end{itemize}
Repeat electrodiagnostic studies are generally not performed unless symptoms persist or recur, suggesting incomplete release or another underlying issue.

\section*{Epidemiology}
Carpal tunnel syndrome (CTS) stands as the most prevalent peripheral nerve entrapment neuropathy, impacting a substantial segment of the adult population globally. Its estimated incidence ranges from 1 to 3 cases per 1,000 persons per year in the general population, with prevalence rates reported between 3\% and 6\%. Women are disproportionately affected compared to men, exhibiting a female-to-male ratio of approximately 3:1 to 5:1, and the peak incidence typically occurs between 40 and 60 years of age, although it can manifest at any age. Common indications for surgical intervention include idiopathic CTS, which is often exacerbated by repetitive hand and wrist movements, as well as CTS associated with systemic conditions such as diabetes mellitus, hypothyroidism, rheumatoid arthritis, and pregnancy. Carpal tunnel release is consistently ranked among the most frequently performed hand surgeries worldwide. While the mortality rate directly attributable to carpal tunnel release is exceedingly low, primarily linked to general anesthesia risks, the morbidity rate is also low. Minor complications, such as scar tenderness or pillar pain, occur in a small percentage of patients, typically resolving over time. Severe complications, including iatrogenic median nerve injury, are rare, estimated at less than 1\% in experienced hands.

\section*{Outcomes \& Prognosis}
Carpal tunnel release is widely regarded as a highly effective surgical intervention, boasting reported success rates for significant symptom relief ranging from 90\% to 95\%. The majority of patients experience substantial improvement in pain, numbness, and tingling, often within days to weeks following the procedure. Functional outcomes, including grip strength, pinch strength, and overall dexterity, typically improve progressively over several months as the median nerve recovers and reinnervates the thenar muscles. The long-term prognosis is generally excellent, with sustained relief for the vast majority of patients. Recurrence of carpal tunnel syndrome after a successful primary release is uncommon, estimated at 1-5\% over a 5-10 year period, and may be attributed to an incomplete initial release, significant re-scarring, or the progression of an underlying systemic condition.

Prognostic factors influencing outcomes include the severity and duration of preoperative symptoms; patients with milder symptoms and a shorter duration of disease generally experience faster and more complete recovery. The presence of thenar muscle atrophy preoperatively may indicate more severe nerve damage and can be associated with less complete motor recovery, although sensory symptoms often improve significantly. Co-morbidities such as diabetes mellitus can also influence recovery rates and the potential for complete symptom resolution. While survival is not a direct outcome measure for this procedure, carpal tunnel release significantly improves the quality of life, functional independence, and ability to perform daily activities and work for affected individuals. Landmark studies and meta-analyses consistently support the efficacy and favorable long-term outcomes of both open and endoscopic carpal tunnel release.

\section{References}
\begin{enumerate}
  \item MedlinePlus. Carpal tunnel release. U.S. National Library of Medicine; 2024. Available from: \url{https://medlineplus.gov/ency/article/002986.htm}
  \item Azar FM, Beaty JH, Canale ST. Campbell's Operative Orthopaedics. 14th ed. Philadelphia, PA: Elsevier; 2021.
  \item American Academy of Orthopaedic Surgeons. Management of Carpal Tunnel Syndrome: Clinical Practice Guideline. Rosemont, IL: AAOS; 2016.
  \item Tintinalli JE, Ma OJ, Yealy DM, et al. Tintinalli's Emergency Medicine: A Comprehensive Study Guide. 9th ed. New York, NY: McGraw-Hill Education; 2020.
\end{enumerate}

\clearpage